\section{Compatibility, descent, and faithful realization}
\label{sec:compatibility}

For a map \(q\colon A\to Q\), a collision is a pair \((a,a')\) with
\(q(a)=q(a')\). The collision space \(A\times_Q A\) is the pullback of the
diagonal of \(Q\) along \((q,q)\). The diagonal of \(A\) is its distinguished
subset, and the residual consists of the distinct pairs with equal images.
Every comparison below specifies the equality being imposed and the objects
on which it is imposed.

\begin{proposition}[Descent along a comparison]\label{prop:descent}
Let \(q\colon A\twoheadrightarrow Q\) be surjective and \(f\colon A\to Z\).
There is a unique \(\bar f\colon Q\to Z\) with \(\bar f q=f\) if and only if
\[
  q(a)=q(a')\quad\Longrightarrow\quad f(a)=f(a').
\]
In particular, \(D\colon A\to A\) descends to an endomorphism of \(Q\)
exactly when \((D,D)\) preserves \(A\times_Q A\).
\end{proposition}

\begin{proof}
If \(\bar f q=f\), apply \(\bar f\) to the equality of images.
Conversely, define \(\bar f(x)=f(a)\) using a representative \(a\) with
\(q(a)=x\). The collision condition makes the value independent of the
representative. Surjectivity gives uniqueness. For the endomorphism
statement, take \(f=qD\).
\end{proof}

A realization \(E\colon V\to Z\) with a left inverse \(J\) reflects
equality. Applying \(J\) to \(E(v)=E(v')\) gives \(v=v'\). If \(V\) is
itself a quotient, this detects equality of its classes. The descent
condition concerns an operation on presentations, while the left inverse
concerns the information retained by a realization.

\subsection{The coding and provability squares}

Fix a G\"odel coding \(g\) with tracked numeral substitution. Let
\(\mathrm{Form}_1\) consist of formulas with at most one free variable,
designated \(x\), and let \(\mathrm{Sent}\) be the sentences. The restrictions
\(g_1\) and \(g_0\) take values in their respective images
\(\mathrm{Code}_1\) and \(\mathrm{Code}_0\). With \(\overline n\) the numeral
for \(n\), define
\[
  \delta_g(\varphi)=\varphi[\overline{g_1(\varphi)}/x],
  \qquad d_g(n)=\operatorname{sub}_g(n,n).
\]
The numeral substitution law gives the exact identity
\(d_g g_1=g_0\delta_g\) \cite{grabmayr,goldfarb}.
\begin{equation}\label{eq:coding}
\begin{CD}
  \mathrm{Form}_1 @>{\delta_g}>> \mathrm{Sent}\\
  @V{g_1}VV @VV{g_0}V\\
  \mathrm{Code}_1 @>{d_g}>> \mathrm{Code}_0 .
\end{CD}
\end{equation}

For a theory \(T\) with sufficient arithmetic and a unary formula
\(\theta(z)\), the diagonal lemma supplies a sentence \(G\) with
\[
  T\vdash G\leftrightarrow\theta(\overline{g_0(G)})
  \qquad\text{\cite{goldfarb,openlogic}}.
\]
Writing \(q_T\colon\mathrm{Sent}\to H_T\) for the quotient by provable
equivalence and \(h_\theta(n)=[\theta(\overline n)]_T\), this is the equality
\(q_T(G)=h_\theta(g_0(G))\) in \(H_T\). It is an agreement at the
constructed sentence \(G\). The coding square \eqref{eq:coding} instead
holds for every formula in its domain. To descend own-code substitution
through an equivalence on formulas, one must also show that equivalent
formula presentations yield equivalent output sentences, as in
Proposition~\ref{prop:descent}.

\subsection{Provability and selected semantics}

For a family \(\mathcal M\) of models of \(T\), let \(Q_{\mathcal M}\)
be the quotient of sentences by equivalence in every model in the family.
Soundness gives a surjective Boolean ring homomorphism
\(c_{\mathcal M}\colon H_T\to Q_{\mathcal M}\), with
\[
 Q_{\mathcal M}\cong H_T/\ker c_{\mathcal M}.
\]
Writing \(I_{\mathcal M}=\ker c_{\mathcal M}\), the map
\((x,y)\mapsto (x,y-x)\) gives a bijection
\[
 H_T\times_{Q_{\mathcal M}}H_T\;\cong\; H_T\times I_{\mathcal M}.
\]
Its inverse adds the difference to the first class, and equal pairs
correspond to the zero element of the ideal.

For a sound, computably axiomatized extension of Robinson arithmetic,
take \(\mathcal M=\{\mathbb N\}\). A true unprovable G\"odel sentence
gives distinct classes \([G]_T\) and \([\top]_T\) with equal semantic image.
In the corresponding ideal, \([\neg G]_T\) is nonzero. When
\(\mathcal M\) contains all models of \(T\), first-order completeness
makes the comparison injective and the residual empty
\cite[Chapters 23 and 37]{openlogic}.
